\chapter{Engenharia de Software}

\section{Produto, processo e qualidade}

Engenharia de software organiza práticas para especificar, construir, testar, entregar, operar e evoluir sistemas com qualidade previsível. Em prova, não reduza a disciplina a ``programar''. Requisitos, arquitetura, testes, configuração, entrega, operação, observabilidade e feedback fazem parte do ciclo de vida.

\begin{teoria}[iterativo e incremental]
\termo{Iterativo} significa revisitar e refinar a solução em ciclos. \termo{Incremental} significa entregar parcelas funcionais crescentes do produto. Um processo pode ser ambos: a cada iteração, aprende-se e melhora-se o produto, entregando um novo incremento potencialmente utilizável.
\end{teoria}

\section{Requisitos e experiência do usuário}

Requisitos funcionais descrevem comportamentos/serviços; não funcionais expressam atributos e restrições como desempenho, segurança, disponibilidade, acessibilidade e auditabilidade. Regras de negócio restringem como a organização opera.

Elicitação pode envolver entrevistas, workshops, observação, análise documental, protótipos e pesquisa com usuários. O gerenciamento de requisitos inclui priorização, rastreabilidade, versionamento, tratamento de mudança e validação.

\begin{exemplo}[história e critério de aceitação]
\textbf{História:} ``Como auditor, quero filtrar contratos por órgão e risco para priorizar fiscalização.''

\textbf{Critério de aceitação:} Dado que existam contratos classificados, quando o auditor selecionar órgão e faixa de risco, então o sistema deve listar apenas registros que satisfaçam ambos os filtros e informar a quantidade total.
\end{exemplo}

Critérios de aceitação devem ser verificáveis. Prototipação reduz incerteza de entendimento, mas um protótipo não é automaticamente arquitetura final nem requisito completo.

\section{Projeto centrado no usuário e storytelling}

Projeto centrado no usuário considera contexto real, necessidades, limitações e ciclos de avaliação. Storytelling ajuda a comunicar jornada, problema e valor do produto, mas não substitui especificação testável quando ela é necessária.

\section{MVP}

Minimum Viable Product é a menor versão que permite testar hipótese relevante de valor/aprendizado com usuários reais. ``Mínimo'' não significa produto sem qualidade essencial, inseguro ou juridicamente inadequado. Em governo, requisitos de segurança, acessibilidade, proteção de dados e conformidade podem ser não negociáveis já no MVP.

\section{Scrum}

Scrum organiza trabalho adaptativo em ciclos. Papéis/accountabilities fundamentais: Product Owner, Scrum Master e Developers, formando o Scrum Team. Eventos incluem Sprint, Sprint Planning, Daily Scrum, Sprint Review e Sprint Retrospective. Artefatos: Product Backlog, Sprint Backlog e Increment; associados a compromissos como Product Goal, Sprint Goal e Definition of Done.

\begin{cebraspe}
A Daily Scrum não é reunião de prestação de contas ao Scrum Master. Seu objetivo é inspecionar progresso em direção ao Sprint Goal e adaptar o Sprint Backlog. O Product Owner maximiza valor e gerencia efetivamente o Product Backlog; não é mero ``escrivão de requisitos''.
\end{cebraspe}

\section{Kanban}

Kanban visualiza fluxo, limita trabalho em progresso (WIP), gerencia fluxo e favorece melhoria contínua. Métricas úteis incluem lead time, cycle time, throughput e WIP. Pela Lei de Little, em sistema estável, uma relação aproximada é:
\[
L=\lambda W
\]
onde $L$ é quantidade média no sistema, $\lambda$ taxa média de chegada/saída e $W$ tempo médio no sistema.

\section{TDD, BDD e ATDD}

\begin{table}[ht]
\centering
\begin{tabularx}{\textwidth}{l X}
\toprule
Prática & Foco \\
\midrule
TDD & ciclo teste que falha $\rightarrow$ código mínimo $\rightarrow$ refatoração; orienta design e verificação automatizada. \\
BDD & comportamento observável e linguagem compartilhada entre negócio e tecnologia; cenários frequentemente em Given/When/Then. \\
ATDD & critérios/testes de aceitação definidos antes da implementação, colaborativamente. \\
\bottomrule
\end{tabularx}
\end{table}

Nenhuma dessas práticas elimina outros testes. TDD não prova ausência de defeitos; BDD não é sinônimo de ferramenta específica.

\section{Testes de software}

\subsection{Níveis e objetivos}

\begin{itemize}
\item \textbf{Unitário}: unidade isolada, rápido e localizado.
\item \textbf{Integração}: interfaces e colaboração entre componentes.
\item \textbf{Funcional/sistema}: comportamento do sistema frente aos requisitos.
\item \textbf{Aceitação}: validação do produto em relação a necessidades/critérios.
\item \textbf{Desempenho}: latência, throughput, recursos e comportamento sob carga.
\item \textbf{Carga}: comportamento em carga esperada/variada.
\item \textbf{Vulnerabilidade/segurança}: busca fragilidades e falhas de controles.
\end{itemize}

Pirâmide de testes é heurística: muitos testes rápidos de baixo nível, menos testes caros ponta a ponta. Contexto pode justificar outras distribuições.

\subsection{Caixa-preta e caixa-branca}

Caixa-preta deriva casos de comportamento especificado: classes de equivalência, valores-limite, tabelas de decisão, transição de estados. Caixa-branca usa estrutura interna: cobertura de instruções, decisões, condições e caminhos.

\begin{atencao}
Cobertura de código mede execução de estruturas, não correção semântica. 100\% de linhas executadas não garante bons asserts, ausência de defeitos ou cobertura de requisitos.
\end{atencao}

\section{Análise estática e SonarQube}

Análise estática inspeciona código/artefatos sem executar o programa. Pode identificar vulnerabilidades, bugs prováveis, duplicação, complexidade e problemas de manutenibilidade. SonarQube agrega regras e métricas e pode atuar em quality gates de CI. Resultado de ferramenta deve ser interpretado; falso positivo e falso negativo existem.

\section{Métricas}

Métricas podem cobrir produto, processo e operação: complexidade ciclomática, cobertura, densidade de defeitos, lead time de mudança, frequência de implantação, MTTR e taxa de falha de mudanças. Métrica isolada pode gerar comportamento disfuncional quando vira meta sem contexto.

\section{DevOps e CI/CD}

Integração contínua incentiva commits frequentes, build automatizado e testes. Entrega contínua mantém software em estado implantável; implantação contínua automatiza promoção para produção quando critérios são atendidos. Pipeline típico:

\begin{center}
\begin{tikzpicture}[node distance=5mm]
\tikzset{s/.style={draw=azulPetroleo,fill=azulClaro,rounded corners,minimum width=2cm,minimum height=.75cm}}
\node[s] (c) {Commit}; \node[s,right=of c] (b) {Build}; \node[s,right=of b] (t) {Testes};
\node[s,right=of t] (sec) {Segurança}; \node[s,right=of sec] (p) {Pacote};
\node[s,below=of p] (d) {Deploy}; \node[s,left=of d] (m) {Monitorar};
\draw[-Latex,thick] (c)--(b)--(t)--(sec)--(p)--(d)--(m);
\end{tikzpicture}
\end{center}

Segregação de funções, aprovação, trilha de auditoria, assinatura de artefatos e gestão de segredos podem ser controles essenciais em pipelines governamentais.

\section{Docker e Kubernetes}

Imagem é artefato imutável em camadas; contêiner é instância em execução. Dockerfile descreve construção. Volume mantém dados fora da camada efêmera do contêiner.

No Kubernetes, \termo{Pod} é unidade básica de execução; \termo{Deployment} gerencia réplicas e atualizações de aplicações stateless; \termo{Service} oferece acesso estável; \termo{ConfigMap} guarda configuração não secreta; \termo{Secret} guarda dados sensíveis, embora proteção efetiva dependa de criptografia, RBAC e controles do cluster; \termo{Ingress/Gateway} expõe tráfego conforme arquitetura.

\begin{cebraspe}
Kubernetes não é mecanismo de backup, não torna aplicação automaticamente stateless e não corrige código não resiliente. Readiness probe responde se instância pode receber tráfego; liveness ajuda a detectar necessidade de reinício; startup probe pode proteger aplicações de inicialização lenta.
\end{cebraspe}

\section{Observabilidade}

Monitoramento acompanha estados conhecidos; observabilidade procura permitir inferir estado interno a partir de sinais. Três sinais clássicos: métricas, logs e traces. Correlação por identificadores facilita seguir requisição entre microsserviços. SLI mede característica; SLO define alvo; SLA é acordo com consequência/compromisso formal.

\section{Mapa mental}

\mapamental{Engenharia de software}{
 child {node[concept]{Requisitos} child {node[concept]{histórias}} child {node[concept]{aceitação}}}
 child {node[concept]{Ágil} child {node[concept]{Scrum}} child {node[concept]{Kanban}}}
 child {node[concept]{Qualidade} child {node[concept]{TDD/BDD}} child {node[concept]{testes}}}
 child {node[concept]{CI/CD} child {node[concept]{pipeline}} child {node[concept]{Sonar}}}
 child {node[concept]{Containers} child {node[concept]{Docker}} child {node[concept]{K8s}}}
 child {node[concept]{Operação} child {node[concept]{observabilidade}}}
}

\section{Questões por nível}

\begin{questao}{\nivel{N1} 12.1}
Qual alternativa descreve corretamente um processo incremental?
\begin{enumerate}[label=\Alph*)]
\item Todo o sistema é especificado e entregue apenas ao final.
\item O produto cresce por parcelas funcionais sucessivas.
\item O código nunca é revisitado.
\item Cada ciclo obrigatoriamente elimina requisitos anteriores.
\item Incremento é sinônimo de protótipo descartável.
\end{enumerate}
\end{questao}
\begin{solucao}{12.1}
\textbf{Gabarito: B.} Incrementalidade adiciona capacidade em parcelas. A descreve abordagem de grande entrega única; C contradiz evolução; D é falsa; E confunde incremento potencialmente útil com protótipo descartável.
\end{solucao}

\begin{questao}{\nivel{N2} 12.2}
Uma equipe possui 100\% de cobertura de linhas, mas defeitos de regra de negócio continuam chegando à produção. A conclusão correta é:
\begin{enumerate}[label=\Alph*)]
\item cobertura de linhas prova ausência de defeitos, logo os relatos são inválidos.
\item cobertura mede execução, não qualidade dos asserts nem suficiência dos cenários.
\item a solução é remover testes unitários.
\item SonarQube substitui testes funcionais.
\item apenas testes de carga detectam regras de negócio incorretas.
\end{enumerate}
\end{questao}
\begin{solucao}{12.2}
\textbf{Gabarito: B.} Cobertura é indicador estrutural. A é absolutista; C elimina feedback útil; D confunde análise estática com validação dinâmica; E restringe indevidamente finalidade de teste de carga.
\end{solucao}

\begin{questao}{\nivel{N3} 12.3}
Em Kubernetes, uma aplicação está viva, porém temporariamente incapaz de atender porque ainda carrega cache obrigatório. O mecanismo mais apropriado para evitar encaminhamento de tráfego durante esse período é:
\begin{enumerate}[label=\Alph*)]
\item readiness probe.
\item apenas liveness probe que sempre falha.
\item apagar o Deployment.
\item aumentar cobertura unitária.
\item usar RAID 5.
\end{enumerate}
\end{questao}
\begin{solucao}{12.3}
\textbf{Gabarito: A.} Readiness informa se o Pod está apto a receber tráfego. Liveness responde à saúde para reinício e não deve ser usada para derrubar processo saudável por indisponibilidade transitória de atendimento. C é desproporcional; D e E não tratam roteamento no cluster.
\end{solucao}

\begin{questao}{\nivel{N4} 12.4}
Auditoria identifica pipeline com deploy automático em produção, credencial administrativa fixa no repositório, ausência de revisão para mudança crítica e logs apagáveis pelo mesmo usuário que executa deploy. Qual combinação representa prioridade de controle mais adequada?
\begin{enumerate}[label=\Alph*)]
\item manter credencial no repositório, mas aumentar tamanho da senha.
\item retirar segredos do código, usar gestão de identidade/segredos, aplicar segregação/aprovação baseada em risco e proteger trilha de auditoria contra alteração indevida.
\item substituir Git por FTP.
\item remover automação e realizar todas as implantações manualmente.
\item apenas aumentar cobertura de testes.
\end{enumerate}
\end{questao}
\begin{solucao}{12.4}
\textbf{Gabarito: B.} O caso envolve segredo exposto, privilégio, segregação e integridade de evidência. A continua expondo segredo; C piora rastreabilidade; D troca risco de automação por erro manual sem atacar governança; E é necessário para qualidade, mas insuficiente para segurança e accountability.
\end{solucao}

\begin{discursiva}[CI/CD e controle]
Descreva, em até 30 linhas, controles mínimos que um Auditor/TI esperaria em pipeline CI/CD de sistema que processa dados pessoais e atos de fiscalização, contemplando código-fonte, testes, segredos, artefatos, aprovação, implantação e logs.
\end{discursiva}
